\documentclass[11pt]{article}
\usepackage[a4paper,margin=2.5cm]{geometry}
\usepackage[table]{xcolor}
\usepackage{titlesec}
\usepackage{enumitem}
\usepackage{longtable}
\usepackage{array}
\usepackage{booktabs}
\usepackage{fancyhdr}
\usepackage{hyperref}
\usepackage{parskip}
\usepackage{ragged2e}
\usepackage{lastpage}
\usepackage{pgfgantt}
\usepackage{mathpazo}
\usepackage[scaled=0.92]{helvet}

\definecolor{primary}{RGB}{0, 75, 141}     % Deep Blue
\definecolor{secondary}{RGB}{204, 85, 0}   % Burnt Orange
\definecolor{lightbg}{RGB}{245, 247, 250}  % Light background

\hypersetup{colorlinks=true, linkcolor=primary, urlcolor=primary}

\titleformat{\section}{\normalfont\Large\bfseries\color{primary}}{\thesection}{0.6em}{}
\titleformat{\subsection}{\normalfont\large\bfseries\color{secondary}}{\thesubsection}{0.6em}{}
\titlespacing*{\section}{0pt}{18pt}{8pt}
\titlespacing*{\subsection}{0pt}{12pt}{5pt}

\pagestyle{fancy}
\fancyhf{}
\renewcommand{\headrulewidth}{0.4pt}
\fancyhead[L]{\small Automated Ranking Portal (ARP)}
\fancyhead[R]{\small DORA}
\fancyfoot[C]{\small\thepage}

\newcolumntype{P}[1]{>{\RaggedRight\arraybackslash}p{#1}}
\setlength{\tabcolsep}{5pt}
\renewcommand{\arraystretch}{1.15}
\setlength{\parindent}{0pt}

\begin{document}

% ---------------- TITLE PAGE ----------------
\begin{titlepage}
\thispagestyle{empty}
\begin{center}
\vspace*{1.0cm}
{\large\bfseries\color{primary} IIT Mandi}\\[4pt]
{\normalsize\color{darkgray} Office of the Dean of Resources \& Alumni Affairs (DORA)}\\[1.8cm]
{\Large\bfseries\color{secondary} Project Proposal and Working Agreement}\\[10pt]
{\huge\bfseries\color{primary} Automated Ranking Portal}\\[6pt]
{\Large\bfseries\color{primary} (ARP)}\\[10pt]
{\large\bfseries An Institutional Ranking Data Aggregation, Processing, and Analytics Platform}\\[1.2cm]
{\color{primary}\rule{0.85\textwidth}{1.5pt}}\\[10pt]
{\normalsize A proposal submitted for consideration and approval by the\\ Office of the Dean of Resources \& Alumni Affairs (DORA), IIT Mandi}\\[10pt]
{\color{primary}\rule{0.85\textwidth}{1.5pt}}
\vspace{1.2cm}

\begin{minipage}[t]{0.48\textwidth}
\textbf{Submitted by:}\\[4pt]
Student Development Team\\
(Team of 4 Students)\\[6pt]
1. \underline{\hspace{4.2cm}}\\[4pt]
2. \underline{\hspace{4.2cm}}\\[4pt]
3. \underline{\hspace{4.2cm}}\\[4pt]
4. \underline{\hspace{4.2cm}}
\end{minipage}
\hfill
\begin{minipage}[t]{0.48\textwidth}
\textbf{Submitted to:}\\[4pt]
Office of the Dean of Resources \& Alumni Affairs (DORA)\\
{}IIT Mandi\\[6pt]
\textbf{Date:} \today
\end{minipage}

\vfill
{\small This document constitutes a proposal and, upon signature by both parties, a working agreement governing scope of work, timelines, reporting, and compensation for the project described herein.}
\end{center}
\end{titlepage}

\tableofcontents
\clearpage

% ---------------- SECTION 1 ----------------
\section{Introduction and Project Overview}
The data collection and compilation process for institutional ranking submissions (such as NIRF, QS World University Rankings, and Times Higher Education) at IIT Mandi is currently a manual pipeline operated by administrative staff. This process involves gathering student, faculty, and learnability data from various departments, manually identifying errors, and reformatting this data into agency-specific formats.

Because data arrives via email in multimodal formats (Excel sheets, Word documents, and PDFs), the process is highly time-consuming and error-prone. This document proposes the design and iterative development of a web-based platform, the \textbf{Automated Ranking Portal} (``ARP''/``the Platform''), by a team of four (4) students (``the Development Team''/``the Team'') under the guidance of DORA.

The Platform covers three connected areas of DORA's work:
\begin{enumerate}[leftmargin=1.6em, itemsep=3pt]
\item \textbf{Data Collection and Governance}: Department nodal officer management, multimodal document ingestion via the portal, secure document repositories for proofs, and automated deadline tracking.
\item \textbf{Processing and Validation}: Automated data unification, AI-assisted error spotting, maker-checker correction workflows, and audit trails for all data modifications.
\item \textbf{Export and Analytics}: Agency-specific template generation (NIRF, QS, THE), executive dashboards for historical tracking, and department-wise contribution/performance analytics.
\end{enumerate}

While the idea is simple to state, the underlying logic involves substantial data schema unification and pipeline complexity (Section~\ref{sec:scope}), so the Platform is best understood as an extensive data governance and compliance engine coupled with a leadership-facing analytics web application. This document also serves as a working agreement defining scope, reporting, review, and compensation for the Development Team.

% ---------------- SECTION 2 ----------------
\section{Scope of Work}
\label{sec:scope}
The scope below reflects the current high-level understanding of requirements. As described in Section~\ref{sec:iterative}, this scope is expected to evolve as DORA raises new requirements during the engagement.

\subsection*{Part A --- Department Data Collection and Governance}
\addcontentsline{toc}{subsection}{Part A --- Department Data Collection and Governance}

\subsection{Institute Organogram and Hierarchy Mapping}
\begin{itemize}[leftmargin=1.4em, itemsep=3pt]
\item Creation of an interactive institute organogram mapping academic, administrative, and functional sections/subsections.
\item Structural representation ensuring data collection is managed according to the exact hierarchical organization of the institute.
\end{itemize}

\subsection{Targeted Admin-Based Data Collection}
\begin{itemize}[leftmargin=1.4em, itemsep=3pt]
\item Admin interface to manage the roster of designated Nodal Officers for each department and section.
\item Capability for the admin to select specific sections based on the institute organogram, identify required data, and dispatch targeted email requests.
\item Automated provisioning of restricted portal access for Nodal Officers to securely upload required ranking proofs.
\end{itemize}

\subsection{Flexible Data Ingestion Engine}
\begin{itemize}[leftmargin=1.4em, itemsep=3pt]
\item Secure submission of all department data exclusively through the portal by authenticated Nodal Officers, eliminating dependency on email-based collection.
\item Submission of required data through Platform-specific web forms and standardised Excel templates, with supporting documents uploaded alongside as evidence.
\item Conversion of multimodal raw data into a single, unified internal schema that maps institutional metrics across different ranking methodologies.
\end{itemize}

\subsection{Secure Document Proof Repository}
\begin{itemize}[leftmargin=1.4em, itemsep=3pt]
\item A centralized, searchable repository to store all raw documents submitted by departments (e.g., scanned proof of research grants, faculty awards) ensuring permanent evidence linkage.
\item Strict data provenance mapping: every single data point must be permanently linked to its source proof (e.g., clicking on a ``Sponsored Research Revenue'' figure allows the auditor to instantly trace its origin and view the exact grant sanction PDF).
\item Document versioning to ensure historical records are maintained when a department submits updated or corrected files.
\end{itemize}

\subsection{Deadline and Compliance Tracking}
\begin{itemize}[leftmargin=1.4em, itemsep=3pt]
\item Admin-defined submission cycles with specified timelines and hard deadlines for data submissions.
\item System tracking of exactly which section has been asked for data, what specific data was requested, expected deadlines, and receipt status.
\item Automated and manual email alerts/reminders dispatched to Nodal Officers whose submissions are pending or incomplete.
\item Compliance dashboard for DORA to view, at a glance, the overall progress of the ranking submission process.
\end{itemize}

\subsection*{Part B --- Data Processing, Unification and Validation}
\addcontentsline{toc}{subsection}{Part B --- Data Processing, Unification and Validation}

\subsection{Historical Data Querying and Value Comparison}
\begin{itemize}[leftmargin=1.4em, itemsep=3pt]
\item Storage and easy backend querying of at least the previous three years of historical ranking data for institutional reference.
\item Real-time contextual display of historically submitted values alongside new data entries during updates (e.g., Previous: X, New: Y, Change: Increase/Decrease).
\item Immediate visual feedback to the data entrant and reviewer, simplifying data verification before submission.
\end{itemize}

\subsection{Faculty and Student Roster Normalization}
\begin{itemize}[leftmargin=1.4em, itemsep=3pt]
\item Standardization of faculty data across departments (e.g., normalizing variations in designation, degree qualifications, and publication indexing details).
\item Aggregation of student demographics, intake numbers, graduation rates, and placement statistics into master institutional metrics.
\item Automated deduplication logic to prevent double-counting of interdisciplinary faculty or joint research projects.
\end{itemize}

\subsection{Automated Error Spotting and Rule-Based Validation}
\begin{itemize}[leftmargin=1.4em, itemsep=3pt]
\item Deterministic rule-based checks applied instantly upon data ingestion (e.g., ensuring `Total PhDs Awarded` cannot be negative, or `Placement Percentage` cannot exceed 100).
\item Automatic flagging of missing mandatory columns, empty rows, or corrupted data cells in uploaded Excel sheets.
\item Visual highlighting of errors within the admin dashboard, allowing the reviewer to quickly trace the error back to the exact source document.
\end{itemize}

\subsection{AI-Assisted Discrepancy Flagging}
\begin{itemize}[leftmargin=1.4em, itemsep=3pt]
\item Heuristic and AI-driven analysis comparing current year submissions against historical trends to detect statistical anomalies.
\item Example: Automatically flagging a department if their reported research publications drop by 80\% compared to the previous year, prompting administrative review.
\item Severity grading of flags (e.g., Critical vs. Warning) to help administrators prioritize data auditing efforts.
\item \textbf{Reliability Principle and Grounding:} The AI layer is an enhancement, never a dependency. The system is designed to work safely with sensible fallbacks if AI services are unavailable. All AI flags must be grounded in approved portal data, avoiding hallucinated policies or numbers.
\end{itemize}

\subsection{Maker-Checker Workflow for Data Corrections}
\begin{itemize}[leftmargin=1.4em, itemsep=3pt]
\item When an error is flagged, the Platform allows DORA to formally raise a ``Correction Request'' directed back to the department's Nodal Officer.
\item The Nodal Officer (Maker) submits the corrected data or overriding justification via the portal.
\item DORA (Checker) reviews the updated submission and explicitly clicks ``Approve'' or ``Reject'', recording the decision. Data is not finalized until Checker approval.
\end{itemize}

\subsection{Audit Log and Data Change Tracking}
\begin{itemize}[leftmargin=1.4em, itemsep=3pt]
\item Immutable recording of every create, edit, approve, reject, and export action across the platform, including the old value, new value, actor, and timestamp, guaranteeing complete data provenance.
\item Dedicated interface for internal auditors to filter, search, and export the audit log to ensure data integrity prior to external submission.
\end{itemize}

\subsection*{Part C --- Agency Export and Ranking Calculation}
\addcontentsline{toc}{subsection}{Part C --- Agency Export and Ranking Calculation}

\subsection{Agency-Specific Schema Mapping}
\begin{itemize}[leftmargin=1.4em, itemsep=3pt]
\item Maintenance of internal configuration maps that translate the unified institutional data into the exact nomenclatures and structures demanded by specific agencies.
\item Support for dynamic parameter adjustments, as ranking agencies frequently update their required data formats year-over-year.
\end{itemize}

\subsection{Final Data Packet Generation and Export}
\begin{itemize}[leftmargin=1.4em, itemsep=3pt]
\item Automated compilation of the approved data into the required final formats (e.g., NIRF Data Capturing System (DCS) Excel templates).
\item Generation of consolidated PDF summary reports for the Director's final sign-off before upload.
\item Supported export targets include \textbf{NIRF} (all 4 categories prioritized), \textbf{MDRA} (Best University \& Engineering), \textbf{QS World University Rankings}, and \textbf{Times Higher Education (THE)}.
\end{itemize}

\subsection*{Part D --- Insights, Dashboards and Analytics}
\addcontentsline{toc}{subsection}{Part D --- Insights, Dashboards and Analytics}

\subsection{Public-Facing Ranking Data \& Visualizations}
\begin{itemize}[leftmargin=1.4em, itemsep=3pt]
\item A publicly accessible portal requiring no administrative or developer access, allowing any visitor to view institutional ranking data.
\item Automated fetching and visualization of year-wise ranking trends across different ranking systems to show performance over time.
\item Visual comparative analytics benchmarking IIT Mandi's performance against other relevant institutions, particularly 2nd-generation IITs.
\end{itemize}

\subsection{Executive Ranking Dashboards}
\begin{itemize}[leftmargin=1.4em, itemsep=3pt]
\item A high-level, leadership-facing dashboard visualizing IIT Mandi's historical ranking performance over a selectable time range (e.g., last 5 years).
\item Parameter-wise breakdown showing exactly where the institute gained or lost points (e.g., Teaching \& Learning vs. Research \& Professional Practice).
\item Visual gap analysis highlighting metrics where the institute is falling behind peer benchmarks or its own historical averages.
\end{itemize}

\subsection{Department-wise Contribution/Performance Analytics}
\begin{itemize}[leftmargin=1.4em, itemsep=3pt]
\item Granular dashboards that analyze internal department contributions to the overall institutional score (e.g., which department brought in the most research funding or had the highest placement rate).
\item Ability for department heads to view their own analytics to identify areas of excellence and areas requiring improvement.
\end{itemize}

\subsection{Automated Briefing Generation}
\begin{itemize}[leftmargin=1.4em, itemsep=3pt]
\item Ability to generate a standardized, boardroom-ready PDF briefing detailing the year's ranking submission statistics, major discrepancies caught by the system, and projected ranking outcomes (which are scenario-based estimates only).
\end{itemize}

\subsection*{Part E --- Cross-Cutting Requirements}
\addcontentsline{toc}{subsection}{Part E --- Cross-Cutting Requirements}

\subsection{Roles and Access Control}
\label{sec:roles}
The Platform will support strict role-based access control (RBAC) as follows:
\begin{itemize}[leftmargin=1.4em, itemsep=3pt]
\item \textbf{Admin/Staff (Data Officer):} Full access to trigger ingestion, review and approve Nodal Officer submissions, initiate correction requests, and generate agency exports.
\item \textbf{Department Nodal Officer:} Restricted access to upload data/proofs for their specific department only, and to respond to assigned correction requests. Cannot view other departments' raw data.
\item \textbf{Leadership/Director:} Read-only access to all executive dashboards, analytics, and final generated reports.
\item \textbf{Auditor:} Read-only access to raw proofs, the audit log, and historical submission trails to ensure compliance.
\end{itemize}

\subsection{Data and System Inputs Required from DORA}
To build and operate the Platform, the Development Team will require the following from DORA:
\begin{itemize}[leftmargin=1.4em, itemsep=3pt]
\item \textbf{Historical ranking data} and previous agency-specific schema formats for NIRF, QS, and THE.
\item \textbf{Sample raw data} (Excel, Word, PDF) typically received from departments to train and test the multimodal ingestion pipeline.
\item \textbf{Authorised email account access:} access to, or a delegated permission on, the authorised email account receiving department submissions.
\item \textbf{Infrastructure support:} guidance on hosting, server, and cloud backup requirements. Where specialized support is needed, DORA will direct the Team to the concerned person/department (e.g., IT Services).
\end{itemize}

\textbf{Confidentiality of institutional data.} All data shared by DORA for the purposes of building and operating the Platform is proprietary and confidential to the Institution. The Development Team shall not share, disclose, publish, or use this data outside the scope of this engagement, and shall not share it with any third party, under any circumstance, without the Institution's prior written consent.

\subsection{Out of Scope (Initial Phase)}
\label{sec:oos}
\begin{itemize}[leftmargin=1.4em, itemsep=3pt]
\item Direct, automated API integration to upload data \textit{into} the NIRF/QS/THE portals, as these agencies typically require manual authenticated uploads by the Director.
\item Historical data entry or digitization of unstructured paper records from previous decades.
\end{itemize}

Any new feature or change requested by DORA beyond the scope above will be evaluated for feasibility. If accepted for inclusion, any change in scope, cost, or timeline shall require written prior approval from DORA before implementation begins.

% ---------------- SECTION 3 ----------------
\section{Development Approach}

\subsection{Iterative Development}
\label{sec:iterative}
A complete, final system architecture cannot be reasonably fixed at the outset given the complexity in Section~\ref{sec:scope}. The Platform will therefore be developed iteratively: new requirements raised by DORA will be assessed for feasibility (Section~\ref{sec:oos}). Any resulting changes to the working scope, project timeline, or compensation shall strictly require prior written approval and will be documented as a signed addendum to this agreement.

\subsection{Version Control and Work Tracking}
All source code will be maintained in an Institute-controlled (DORA-owned) GitHub repository from the outset. All work will be tracked and submitted via GitHub Pull Requests (PRs), which serve as the primary record of individual contribution. The platform and its data will be hosted exclusively on Institute-provisioned servers.

\subsection{Institutional Deployment and Production Readiness}
In line with best practices for institutional systems, the Platform will pass through stringent production gates before real data is loaded:
\begin{itemize}[leftmargin=1.4em, itemsep=3pt]
\item \textbf{Ownership and Control:} IIT Mandi retains full control of the production database, storage, domains, deployment accounts, and the underlying data schema from day one.
\item \textbf{Reproducible Handover:} The team will deliver installation, configuration, backup/restore, and deployment runbooks, ensuring an Institute technical owner can independently manage staging and production instances.
\item \textbf{Controlled Pilot and Go/No-Go:} Before a campus-wide launch, the system will undergo a controlled pilot with selected departments to validate KPIs, followed by a formal go/no-go decision. The pilot must strictly meet the following \textbf{Measurable KPIs}:
  \begin{itemize}[leftmargin=1.4em, itemsep=2pt]
    \item \textit{AI Reliability:} $>$95\% accuracy in flagging data discrepancies, with 0\% false positives on critical institutional metrics.
    \item \textit{Data Auditability:} 100\% of data modifications must be accurately tied to a recorded user session and timestamp in the immutable audit log.
    \item \textit{System Resilience:} A witnessed disaster recovery drill must confirm the target 24-hour RPO and 8-hour RTO are achievable.
  \end{itemize}
\end{itemize}

% ---------------- SECTION 4 ----------------
\section{Team Structure}
The Development Team shall consist of four (4) students, listed below (to be finalized and appended to this document upon submission):

\begin{enumerate}[leftmargin=1.8em, itemsep=4pt]
\item \underline{\hspace{7cm}} (Team Lead / Point of Contact)
\item \underline{\hspace{7cm}} (Team Lead / Point of Contact)
\item \underline{\hspace{7cm}}
\item \underline{\hspace{7cm}}
\end{enumerate}

Team members 1 and 2 shall jointly serve as Team Leads and primary points of contact with DORA for scheduling meetings, consolidating weekly reports, and coordinating delivery.

% ---------------- SECTION 5 ----------------
\section{Timeline and Prioritized Milestones}
\label{sec:timeline}
Based on administrative requirements, the development pipeline is strictly prioritized to deliver maximum immediate value. The timeline is divided into four distinct phases, targeting an overall completion window ending \textbf{14 December 2026} (Section~\ref{sec:duration}).

\subsection{Phase 1: Core NIRF Workflow (High Priority)}
\begin{itemize}[leftmargin=1.4em, itemsep=3pt]
\item \textbf{Deliverable:} A functional end-to-end system for the NIRF data collection, validation, approval, and export lifecycle.
\item \textbf{Scope:} Institute organogram setup, admin targeted data requests, node-officer upload portal, maker-checker validation workflow, and automated generation of the NIRF DCS Excel exports.
\item \textbf{Acceptance Criteria:} The system successfully generates a complete, compliant NIRF export based on mock data inputs, and the maker-checker audit trail correctly logs all data modifications.
\item \textbf{Estimated Effort:} 80 person-hours per week (20 hours per member).
\end{itemize}

\subsection{Phase 2: Historical Analytics \& Multi-Ranking Export}
\begin{itemize}[leftmargin=1.4em, itemsep=3pt]
\item \textbf{Deliverable:} Executive analytics dashboards and integration for additional ranking agencies.
\item \textbf{Scope:} QS, THE, and MDRA schema mapping and exports. Department-wise contribution/performance analytics, and automated boardroom-briefing generation.
\item \textbf{Acceptance Criteria:} Leadership dashboards accurately visualize historical trends and gap analyses. The platform can export valid templates for at least one additional agency (e.g., QS).
\item \textbf{Estimated Effort:} 80 person-hours per week (20 hours per member).
\end{itemize}

\subsection{Phase 3: Public Visualizations \& System Hardening}
\begin{itemize}[leftmargin=1.4em, itemsep=3pt]
\item \textbf{Deliverable:} Public-facing visualization portal and advanced system compliance.
\item \textbf{Scope:} Public ranking trend visualizations, competitor benchmarking. Full implementation of automated reminders, contextual historical-value comparisons during data entry, and comprehensive RBAC testing.
\item \textbf{Acceptance Criteria:} Public portal is accessible without authentication and accurately plots dummy historical data. Automated reminders trigger appropriately based on predefined deadlines.
\item \textbf{Estimated Effort:} 80 person-hours per week (20 hours per member).
\end{itemize}

\subsection{Phase 4: UAT, Handover, and Post-Deployment Support}
\begin{itemize}[leftmargin=1.4em, itemsep=3pt]
\item \textbf{Deliverable:} Production-ready deployment, complete documentation, and trained staff.
\item \textbf{Scope:} Detailed User Acceptance Testing (UAT) with actual DORA staff, comprehensive technical and user manual documentation, training sessions for nodal officers, and a defined post-deployment support period. \textbf{Addition of Controlled Pilot:} A closed pilot with a selected subset of departments, allowing KPI measurement before a full migration and go/no-go decision.
\item \textbf{Acceptance Criteria:} DORA staff formally signs off on UAT. Handover of the Institute-controlled repository is complete. User manuals are provided. The service is functionally operable without dependency on the Development Team.
\item \textbf{Estimated Effort:} 70 person-hours per week (17.5 hours per member).
\end{itemize}

\begin{center}
\begin{ganttchart}[
  vgrid, hgrid,
  x unit=0.42cm,
  y unit title=0.7cm,
  y unit chart=0.62cm,
  title height=1,
  bar height=0.5,
  bar label font=\scriptsize,
  title label font=\scriptsize\color{primary},
  milestone label font=\scriptsize,
  group label font=\scriptsize\bfseries\color{secondary},
  bar/.append style={fill=primary!80, draw=primary, rounded corners=2pt},
  bar label node/.append style={align=right},
  ]{1}{12}
  \gantttitle{Prioritized Project Timeline: Weeks}{12} \\
  \gantttitlelist{1,...,12}{1} \\
  \ganttgroup{Phase 1: Core NIRF Workflow}{1}{4} \\
  \ganttbar{Organogram \& Data Collection}{1}{2} \\
  \ganttbar{Maker-Checker \& NIRF Export}{3}{4} \\
  \ganttgroup{Phase 2: Analytics \& Ext.}{5}{7} \\
  \ganttbar{QS, THE, MDRA Schema Map}{5}{6} \\
  \ganttbar{Executive \& Dept Analytics}{7}{7} \\
  \ganttgroup{Phase 3: Public \& Hardening}{8}{10} \\
  \ganttbar{Public Visualizations}{8}{9} \\
  \ganttbar{Reminders \& History Query}{10}{10} \\
  \ganttgroup{Phase 4: UAT \& Handover}{11}{12} \\
  \ganttbar{Security Audit \& UAT}{11}{11} \\
  \ganttbar{Documentation \& Training}{12}{12} \\
  \ganttbar{Production Go-Live}{12}{12}
\end{ganttchart}
\end{center}

\textit{Note: this Gantt chart illustrates the sequence of the prioritized phases. Actual phase durations, dependencies, and start dates will be agreed with DORA and revised as scope evolves under Section~\ref{sec:oos}. Please see Appendix A for a detailed Milestone Matrix.}

% ---------------- SECTION 6 ----------------
\section{Reporting, Review, and Communication}

\subsection{Weekly Progress Reports}
The Development Team shall submit one (1) written progress report per week, summarizing work completed, work in progress, and blockers, attributed by individual member and referencing the corresponding GitHub Pull Requests (Section 3.2).

\subsection{Weekly Demonstration Meeting}
The Development Team and a designated Administrative representative shall meet at most once per week to demonstrate progress and gather inputs. A meeting is not mandatory every single week; it will be scheduled as needed, by mutual convenience, particularly if there is no material update since the previous meeting.

\subsection{Administrative Review and Delegation}
Final data exports are generated only after review and approval by DORA.

\subsection{Documentation of Communication}
All key communication and decisions exchanged between the Team and DORA, including outcomes of meetings, will additionally be documented over email, serving as minutes of meeting (MoM) for record-keeping.

% ---------------- SECTION 7 ----------------
\section{Compensation and Payment Terms}
\label{sec:compensation}

\subsection{Stipend}
Each member of the Development Team shall be compensated at a rate of Rs.~200 per hour of documented, verifiable work. To ensure budget predictability, the following weekly compensation ceiling is established (based on a maximum of 20 billable hours per member):

\begin{table}[h]
\centering
\renewcommand{\arraystretch}{1.4}
\begin{tabular}{p{5cm} p{4cm} p{4cm}}
\rowcolor{primary}
\textcolor{white}{\textbf{Resource}} & \textcolor{white}{\textbf{Max Hours / Week}} & \textcolor{white}{\textbf{Max Budget / Week}} \\
\rowcolor{lightbg}
Per Team Member & 20 Hours & Rs.~4,000 \\
\rowcolor{white}
\textbf{Total Team (4 Members)} & \textbf{80 Hours} & \textbf{Rs.~16,000} \\
\end{tabular}
\end{table}

Any hours exceeding this ceiling must receive prior written approval from DORA.

\subsection{Estimated Cumulative Cost}
Based on the phase-wise estimated efforts and the 12-week project timeline, the projected cumulative cost of the Development Team is calculated comprehensively as follows:

\begin{table}[h]
\centering
\renewcommand{\arraystretch}{1.4}
\begin{tabular}{p{1.5cm} p{2cm} p{6cm} p{2cm} p{2cm}}
\rowcolor{primary}
\textcolor{white}{\textbf{Phase}} & \textcolor{white}{\textbf{Duration}} & \textcolor{white}{\textbf{Calculation (Members $\times$ Hrs/Wk $\times$ Wks)}} & \textcolor{white}{\textbf{Total Hrs}} & \textcolor{white}{\textbf{Est. Cost}} \\
\rowcolor{lightbg}
Phase 1 & 4 Weeks & 4 Members $\times$ 20 Hrs/Wk $\times$ 4 Wks & 320 Hours & Rs.~64,000 \\
\rowcolor{white}
Phase 2 & 3 Weeks & 4 Members $\times$ 20 Hrs/Wk $\times$ 3 Wks & 240 Hours & Rs.~48,000 \\
\rowcolor{lightbg}
Phase 3 & 3 Weeks & 4 Members $\times$ 20 Hrs/Wk $\times$ 3 Wks & 240 Hours & Rs.~48,000 \\
\rowcolor{white}
Phase 4 & 2 Weeks & 4 Members $\times$ 17.5 Hrs/Wk $\times$ 2 Wks & 140 Hours & Rs.~28,000 \\
\rowcolor{secondary}
\textcolor{white}{\textbf{Total}} & \textcolor{white}{\textbf{12 Wks}} & \textbf{Calculated at a rate of Rs.~200/hr} & \textcolor{white}{\textbf{940 Hrs}} & \textcolor{white}{\textbf{Rs.~188,000}} \\
\end{tabular}
\end{table}

\textit{Detailed Billable Hours Breakdown:}
\begin{itemize}[leftmargin=1.4em, itemsep=3pt]
\item \textbf{Phase 1 (Weeks 1--4):} Core NIRF Workflow requires maximum intensity. 4 members working up to the 20 hours/week ceiling over 4 weeks equates to $4 \times 20 \times 4 = 320$ hours. At Rs.~200/hr, the cost is Rs.~64,000.
\item \textbf{Phase 2 (Weeks 5--7):} Analytics \& Extensions continue at peak effort. 4 members $\times$ 20 hours/week $\times$ 3 weeks equates to $4 \times 20 \times 3 = 240$ hours. Cost is Rs.~48,000.
\item \textbf{Phase 3 (Weeks 8--10):} Hardening phase continues with high intensity. 4 members $\times$ 20 hours/week $\times$ 3 weeks equates to $4 \times 20 \times 3 = 240$ hours. Cost is Rs.~48,000.
\item \textbf{Phase 4 (Weeks 11--12):} UAT and Handover require slightly reduced engagement. 4 members $\times$ 17.5 hours/week $\times$ 2 weeks equates to $4 \times 17.5 \times 2 = 140$ hours. Cost is Rs.~28,000.
\end{itemize}

\textit{Note: This cumulative cost is an estimate based on projected efforts. Payment is strictly based on actual verifiable hours worked, capped at the absolute maximum project ceiling of 960 hours (Rs.~1,92,000) for the 12-week period.}

\subsection{Development Tool Subscription}

To support the accelerated development schedule and parallel work of the four-member Development Team, the following development-support resource is proposed:

\begin{itemize}[leftmargin=1.4em, itemsep=3pt]

\item \textbf{Purpose:} AI-assisted support for repository-level development, debugging, code review, technical investigation, and documentation during the intensive development period.

\item \textbf{Proposed Tool:} Claude Pro, including Claude Code and project-level context capabilities relevant to the Platform's multi-component development workflow.

\item \textbf{Individual Access \& Rationale:} One subscription per team member is required to support concurrent development across frontend, backend, database, testing, and deployment workstreams. Anthropic's legal Terms of Service strictly prohibit sharing a single individual account across multiple users. Furthermore, building a robust, scalable system of this complexity requires processing thousands of lines of code and maintaining large context windows. A single shared account would immediately hit strict rate limits, severely bottlenecking the entire team. Individual subscriptions ensure parallel development streams remain unblocked, allowing the team to deliver a high-quality, production-ready system within the aggressive 12-week timeline.

\item \textbf{Duration and Cost:} Two months for four members at the prevailing individual monthly rate of USD~20 per user, resulting in an estimated total of USD~160, excluding applicable taxes or local pricing adjustments.

\item \textbf{Approval and Usage:} Procurement or reimbursement shall require prior written approval from DORA and shall remain subject to applicable Institute rules. Institutional or personal confidential data shall not be submitted to external AI services without explicit authorization and appropriate safeguards.

\end{itemize}

\subsection{Total Cost and Comparison with Third-Party Alternatives}
The total estimated development cost for the in-house student team is \textbf{Rs.~2,01,500} (comprising the Rs.~1,88,000 maximum stipend and approx. Rs.~13,500 / USD~160 for Claude Pro subscriptions).

If a similar bespoke, data-sensitive enterprise platform were to be commissioned from a third-party software development agency, the development cost (including dedicated UI/UX design, full-stack engineering, cloud architecture, and security testing) would typically range between \textbf{Rs.~10,00,000 to Rs.~25,00,000} or more, often accompanied by exorbitant recurring Annual Maintenance Contracts (AMC). By utilizing the technical expertise of our own students, the Institute is acquiring a highly customized, secure solution for \textbf{less than 20\% of the market cost}, while retaining complete data sovereignty and fostering in-house talent.

% \subsection{Non-Monetary Compensation}
% In addition to the hourly stipend, any development tool or software subscription (e.g., Claude Code access) shall be procured or reimbursed only with prior written approval and strictly as per Institute rules.

\subsection{Payment Schedule}
Compensation shall be calculated and disbursed on a monthly basis. Payment for hours worked in a given calendar month shall be released during the first week of the following month, subject to submission and administrative acceptance of that month's weekly reports.

\subsection{Integrity of Hours and Basis for Payment}
Payment in a given month is contingent on submission of that month's weekly reports and verifiable work reflected through GitHub activity (Section 3.2). The Development Team commits to logging hours with utmost sincerity and honesty; no inflated, false, or unworked hours shall be claimed for compensation.

% ---------------- SECTION 8 ----------------
\section{Intellectual Property and Ownership}
\begin{itemize}[leftmargin=1.4em, itemsep=3pt]
\item All source code, documentation, designs, data, and other work products created or populated by the Development Team under this engagement shall be the exclusive property of the Institution. The source code repository must be Institute-controlled at all times. This includes full institutional ownership of the underlying data and codebase, not merely a license to use the Platform.
\item The Development Team shall be permitted to reference this project as part of their academic or professional portfolio, subject to prior written consent of DORA regarding the extent of disclosure, and excluding any confidential data.
\end{itemize}

% ---------------- SECTION 9 ----------------
\section{Information Security, Data Privacy, and Confidentiality}
\begin{itemize}[leftmargin=1.4em, itemsep=3pt]
\item \textbf{Information Security \& RBAC:} The Platform shall implement strict Role-Based Access Control (RBAC) as defined in Section~\ref{sec:roles}. Security measures against unauthorized access, SQL injection, and cross-site scripting must be incorporated.
\item \textbf{Backup and Disaster Recovery:} Automated daily backups of the database and document repository must be implemented. A disaster recovery plan shall be formulated to restore services with minimal downtime in case of critical failure.
\item The Platform will process sensitive data related to institutional, faculty, and student records. Access to production data shall be limited to what is strictly necessary for development, testing, and debugging, and shall follow any data protection guidelines issued by the Institution.
\item The Development Team shall not use, disclose, or retain institutional data for any purpose outside the scope of this engagement, both during and after the engagement period.
\item \textbf{AI Privacy, Grounding, and Auditability:} As the Platform utilizes AI for discrepancy flagging, answers and flags must be grounded in approved institutional data. The model must operate within a strict data boundary, ensuring sensitive personal records are not exposed to external LLMs without contractual safeguards. All AI requests must be logged for auditability, with clear routing for human escalation.
\item \textbf{Data Minimization and Transparency:} The platform must collect only data required for the approved workflow and maintain a data inventory describing each field, its purpose, retention period, and deletion action.
\end{itemize}

% ---------------- SECTION 10 ----------------
\section{Duration and Termination}
\label{sec:duration}
\begin{itemize}[leftmargin=1.4em, itemsep=3pt]
\item This agreement shall commence on the date of signature by both parties. The Team targets completion and handover of the Platform to DORA by \textbf{14 December 2026} (a 12-week period starting 22 September 2026). This target date is subject to change based on more or fewer requirements being added to scope over the course of the engagement (Section~\ref{sec:oos}); the Development Team will make its best effort to complete the project as soon as possible.
\item In the event of termination, compensation shall be payable for all verified work completed up to the effective date of termination, following Section~\ref{sec:compensation}.
\end{itemize}

% ---------------- SECTION 11 ----------------
\section{Post-Deployment Maintenance and Long-Term Sustainability}
\label{sec:sustainability}

\subsection{Responsibility Allocation Matrix}
Once the 12-week development period concludes and the Platform is handed over in production, the ongoing maintenance and governance must be distributed according to the following institutional owners:

\begin{table}[h]
\centering
\renewcommand{\arraystretch}{1.4}
\begin{tabular}{p{4.5cm} p{10cm}}
\rowcolor{primary}
\textcolor{white}{\textbf{Domain}} & \textcolor{white}{\textbf{Accountable Institutional Owner}} \\
\rowcolor{lightbg}
\textbf{Business Policy \& Content} & \textbf{DORA:} Responsible for deciding NIRF/QS reporting logic, nodal officer rosters, and resolving flagged discrepancies. \\
\rowcolor{white}
\textbf{Platform Operations} & \textbf{IIT Mandi IT Services (or designated technical unit):} Responsible for managing server hosting, domain names, periodic data backups, and technical incident response. \\
\rowcolor{lightbg}
\textbf{Data Protection \& Security} & \textbf{IIT Mandi Security Authority:} Responsible for annual penetration testing, access reviews, and compliance governance. \\
\end{tabular}
\end{table}

\subsection{Recurring Costs and Sustainability}
The estimated budget in Section~\ref{sec:compensation} covers the \textit{one-time development cost}. To ensure the Platform's long-term sustainability (Years 1--3 and beyond), DORA and the Finance office must provision an annual operating budget to cover the following recurring infrastructure costs:
\begin{itemize}[leftmargin=1.4em, itemsep=3pt]
\item Managed relational database hosting and AWS S3 (or equivalent) secure object storage.
\item Institutional domain names, TLS/SSL certificates, and transactional email/SMS service quotas.
\item API token usage costs for the AI-Assisted Discrepancy Flagging models (based on a finalized cost ceiling).
\end{itemize}

% ---------------- SECTION 12 ----------------
\section{Acceptance and Signatures}
By signing below, both parties acknowledge that they have read, understood, and agree to the terms described in this document, including the understanding that the technical scope in Section~\ref{sec:scope} is indicative and will evolve iteratively as described in Section~\ref{sec:iterative}.

\vspace{1cm}
\begin{minipage}[t]{0.48\textwidth}
\textbf{For the Office of the Dean of Resources \& Alumni Affairs (DORA), IIT Mandi}\\[1.4cm]
\underline{\hspace{6.5cm}}\\
Name:\underline{\hspace{5cm}}\\[6pt]
Designation:\underline{\hspace{4.5cm}}\\[6pt]
Date:\underline{\hspace{5.5cm}}
\end{minipage}
\hfill
\begin{minipage}[t]{0.48\textwidth}
\textbf{For the Student Development Team}\\[6pt]
\underline{\hspace{6.5cm}}\\
Name:\underline{\hspace{5cm}}\\
Designation: Team Lead\\
Date:\underline{\hspace{5.5cm}}\\[10pt]
\underline{\hspace{6.5cm}}\\
Name:\underline{\hspace{5cm}}\\
Designation: Team Lead\\
Date:\underline{\hspace{5.5cm}}
\end{minipage}

\vspace{1cm}
\textbf{Other Team Members (Acknowledgement of Terms):}

\begin{tabular}{@{}p{7.5cm}p{6cm}@{}}
1. \underline{\hspace{6cm}} & Signature: \underline{\hspace{3.5cm}} \\[10pt]
2. \underline{\hspace{6cm}} & Signature: \underline{\hspace{3.5cm}} \\
\end{tabular}

\clearpage
\appendix
\section{Appendix A: Milestone and Deliverable Matrix}

\begin{table}[h]
\centering
\renewcommand{\arraystretch}{1.4}
\begin{tabular}{p{2cm} p{3.5cm} p{4.5cm} p{1.8cm} p{2.2cm}}
\rowcolor{primary}
\textcolor{white}{\textbf{Phase}} & \textcolor{white}{\textbf{Key Deliverables}} & \textcolor{white}{\textbf{Acceptance Criteria}} & \textcolor{white}{\textbf{Effort}} & \textcolor{white}{\textbf{Timeline}} \\
\rowcolor{lightbg}
\textbf{Phase 1}: Core NIRF Workflow & End-to-end NIRF data collection, validation, and export platform & System generates a valid NIRF Excel export from mock inputs; maker-checker workflow correctly records audit logs. & 80 hrs/week & Weeks 1--4 \\
\rowcolor{white}
\textbf{Phase 2}: Analytics & Executive dashboards and QS/THE schema mappings & Dashboards correctly visualize historical data; export works for at least one additional ranking agency. & 80 hrs/week & Weeks 5--7 \\
\rowcolor{lightbg}
\textbf{Phase 3}: Public \& Hardening & Public ranking visualizations, historical query, and reminders & Public portal accessible without login; automated email reminders trigger on deadlines. & 80 hrs/week & Weeks 8--10 \\
\rowcolor{white}
\textbf{Phase 4}: UAT \& Handover & Production deployment, technical/user documentation, and training & DORA signs off on UAT; GitHub repo handed over; staff training sessions completed successfully. & 70 hrs/week & Weeks 11--12 \\
\end{tabular}
\caption{One-Page Milestone and Deliverable Matrix}
\end{table}

\vspace{1cm}
\section{Appendix B: Proposed System Architecture}

\textbf{Frontend (Client Layer):}
\begin{itemize}[leftmargin=1.4em, itemsep=3pt]
\item Built using Next.js/React.js to provide a highly responsive, dynamic user experience across desktop and mobile form factors.
\item Implementation of a strict Role-Based Access Control (RBAC) routing mechanism to securely isolate public views, authenticated Nodal Officer portals, and privileged Admin/Executive dashboards.
\item Utilization of institutional navy-and-saffron branding and consistent UX patterns (e.g., empty-state guidance, loading states, and multi-step form confirmations).
\end{itemize}

\textbf{Backend (Application Layer):}
\begin{itemize}[leftmargin=1.4em, itemsep=3pt]
\item Next.js API Routes (Node.js) architected for idempotency, ensuring secure handling of data ingestion, maker-checker validation workflows, and dynamic multi-agency schema mapping logic.
\item Resilient document parsing engine built to process multimodal submissions (Excel, PDF, Word) safely in isolated environments to mitigate malware risks.
\item Integration with Institute's unified authentication and Single Sign-On (SSO) systems, enforcing Multi-Factor Authentication (MFA) for all privileged administrative accounts.
\end{itemize}

\textbf{Data, Storage, and Security Layer:}
\begin{itemize}[leftmargin=1.4em, itemsep=3pt]
\item A managed Relational Database (e.g., PostgreSQL) designed with a robust schema-migration process. It stores structured ranking metrics, user rosters, and append-only immutable audit logs to guarantee strict data provenance.
\item Encrypted Secure Object Storage (e.g., AWS S3 or an Institute-provisioned equivalent) enforcing restricted access URLs with expiring tokens for the centralized document proof repository.
\item End-to-end security measures including TLS for data-in-transit, encryption-at-rest for sensitive database fields, and isolated environment management for credentials.
\item A defined Disaster Recovery Plan (DRP) including automated, encrypted daily backups designed to meet a target Recovery Point Objective (RPO) of 24 hours and a Recovery Time Objective (RTO) of 8 hours, subject to Institute approval.
\end{itemize}

\textbf{Deployment, Hosting, and CI/CD Operations:}
\begin{itemize}[leftmargin=1.4em, itemsep=3pt]
\item Production deployment on Institute-provisioned servers maintaining absolute data sovereignty. Network access will be restricted to designated campus IPs where mandated.
\item A fully Institute-controlled GitHub repository serving as the single source of truth for version control, issue tracking, and architectural documentation.
\item An automated CI/CD pipeline bridging separate testing, staging, and production environments. Every pull request will trigger automated dependency/vulnerability scanning (SAST/DAST) before merge.
\item Implementation of structured operational runbooks specifying support hours, severity levels, escalation contacts, and predefined release rollback procedures.
\end{itemize}

\end{document}
  